\section{Find the First and Last Occurrence of an Element}
\label{sec:bs-first-last}

\problemstatement{We are given a sorted array $\textit{arr}$ of $n$
elements, possibly with duplicates, and a target $X$. We must find the
index of the \textbf{first} occurrence of $X$ and the index of the
\textbf{last} occurrence of $X$. If $X$ is not present, both should be
$-1$.}

\begin{patternbox}
\emph{``First occurrence''} and \emph{``last occurrence''} of a value in a
sorted array are, respectively, exactly the lower bound and
(upper bound $- 1$) of that value, \emph{provided the value is actually
present}. This section is mostly about correctly handling the ``what if it
is not present at all'' case, since lower bound and upper bound alone do
not distinguish between ``$X$ is present, starting here'' and ``$X$ is
absent, and this is just where it would be inserted.''
\end{patternbox}

\subsection*{Understanding the problem carefully}

Recall from Section~\ref{sec:bs-upper-bound} that, when $X$ is present,
\[
  \textit{firstOccurrence} = \textit{lowerBound}(X), \qquad
  \textit{lastOccurrence} = \textit{upperBound}(X) - 1.
\]
The subtlety is verifying presence. The lower bound of $X$ gives the first
index whose value is $\ge X$; this is the same as the first occurrence of
$X$ \emph{only if} the value actually stored at that index equals $X$ (and
the index is within bounds). If $\textit{lowerBound}(X) = n$, or if
$\textit{arr}[\textit{lowerBound}(X)] \ne X$, then $X$ is simply not in the
array at all.

\begin{ideabox}[Approach]
Compute $\textit{lb} \gets \textit{lowerBound}(X)$. If $\textit{lb} = n$ or
$\textit{arr}[\textit{lb}] \ne X$, return $(-1,-1)$. Otherwise, the first
occurrence is $\textit{lb}$, and the last occurrence is
$\textit{upperBound}(X) - 1$.
\end{ideabox}

\subsection*{Why checking presence only once is sufficient}

Because the array is sorted, every occurrence of $X$ forms a single
contiguous block. If the very first candidate position (the lower bound)
does not actually hold $X$, no later position can hold $X$ either -- the
lower bound is, by construction, the leftmost position any element $\ge X$
could occupy, so if even that position fails to be exactly $X$, $X$ does
not appear anywhere in the array. This means a single equality check after
computing the lower bound is enough to settle presence, with no need for
an additional independent search.

\subsection*{Algorithm}

\begin{enumerate}
  \item $\textit{lb} \gets \textit{lowerBound}(\textit{arr}, X)$ (using
        Algorithm~\ref{sec:bs-lower-bound}).
  \item If $\textit{lb} = n$ or $\textit{arr}[\textit{lb}] \ne X$: return
        $(-1, -1)$.
  \item Otherwise: $\textit{ub} \gets \textit{upperBound}(\textit{arr}, X)$
        (using Algorithm~\ref{sec:bs-upper-bound}). Return
        $(\textit{lb}, \textit{ub} - 1)$.
\end{enumerate}

\begin{algorithm}[H]
\caption{First and Last Occurrence}
\begin{algorithmic}[1]
\Procedure{FirstAndLast}{$\textit{arr}, X$}
  \State $\textit{lb} \gets \Call{LowerBound}{\textit{arr}, X}$
  \If{$\textit{lb} = n$ \textbf{or} $\textit{arr}[\textit{lb}] \ne X$}
    \State \Return $(-1, -1)$
  \EndIf
  \State $\textit{ub} \gets \Call{UpperBound}{\textit{arr}, X}$
  \State \Return $(\textit{lb}, \textit{ub} - 1)$
\EndProcedure
\end{algorithmic}
\end{algorithm}

\subsection*{Dry run}

\dryrun{Take $\textit{arr} = [5, 7, 7, 8, 8, 10]$, $X = 8$.}

$\textit{lowerBound}(8)$: scanning, the first index with value $\ge 8$ is
index $3$ ($\textit{arr}[3]=8$). Since $\textit{arr}[3]=8=X$, $X$ is
present. $\textit{upperBound}(8)$: the first index with value $>8$ is
index $5$ ($\textit{arr}[5]=10$). So
$\textit{lastOccurrence} = 5 - 1 = 4$.

The answer is (first $=3$, last $=4$).

Now take $X = 6$ on the same array. $\textit{lowerBound}(6) = 1$
($\textit{arr}[1]=7 \ge 6$ is the first such index). But
$\textit{arr}[1] = 7 \ne 6$, so $6$ is not present, and the answer is
$(-1, -1)$.

\subsection*{C++ implementation}

\begin{lstlisting}
#include <bits/stdc++.h>
using namespace std;

int lowerBound(vector<int>& arr, int X) {
    int n = arr.size();
    int lo = 0, hi = n - 1, ans = n;
    while (lo <= hi) {
        int mid = lo + (hi - lo) / 2;
        if (arr[mid] >= X) { ans = mid; hi = mid - 1; }
        else lo = mid + 1;
    }
    return ans;
}

int upperBound(vector<int>& arr, int X) {
    int n = arr.size();
    int lo = 0, hi = n - 1, ans = n;
    while (lo <= hi) {
        int mid = lo + (hi - lo) / 2;
        if (arr[mid] > X) { ans = mid; hi = mid - 1; }
        else lo = mid + 1;
    }
    return ans;
}

pair<int,int> firstAndLast(vector<int>& arr, int X) {
    int n = arr.size();
    int lb = lowerBound(arr, X);
    if (lb == n || arr[lb] != X) return {-1, -1};

    int ub = upperBound(arr, X);
    return {lb, ub - 1};
}

int main() {
    vector<int> arr = {5, 7, 7, 8, 8, 10};
    auto [first, last] = firstAndLast(arr, 8);
    cout << "First: " << first << ", Last: " << last << endl;
    return 0;
}
\end{lstlisting}

\begin{complexitybox}
\textbf{Time Complexity.} Two binary searches (lower bound and, if needed,
upper bound), each $O(\log n)$, giving an overall time complexity of
\[
  O(\log n).
\]
\textbf{Space Complexity.} $O(1)$ auxiliary space.
\end{complexitybox}

\begin{takeawaybox}
Lower bound and upper bound, used together with a single presence check,
solve any ``range of occurrences'' question about a sorted array. There is
rarely a need to write a bespoke binary search for ``first occurrence'' or
``last occurrence'' from scratch -- compose them from the lower/upper bound
primitives instead.
\end{takeawaybox}
